% Final strengthening of special_row_sieve_theorems.tex.
\section{Restricting the sieve to order ideals of the implication poset}
\label{sec:ideal-sieve}

There is further structure in the possible color masks: ordinary
reachability forces them to be upward closed.  The full Boolean cube
therefore contains states that no spanning tree can induce.  Removing
these states improves both the parameter and the practical table size.

Collapse ordinary SCCs meeting the exceptional supports to colors, and
discard the colors already reachable from $Z$.  Let $\mathcal P$ be the
partial order on the remaining $d$ colors in which
\[
                  u\preceq v
           \quad\Longleftrightarrow\quad
       \text{an ordinary directed path leads from color $u$ to color $v$}.
\]
Distinct SCCs make this relation antisymmetric.  Reachability is taken in
the full ordinary graph, including intermediate rows outside the special
set.  Let $\mathcal I(\mathcal P)$ be the set of order ideals, and write
\[
                        J=|\mathcal I(\mathcal P)|.
\]
Here an ideal $K$ satisfies $v\in K$ and $u\preceq v\Rightarrow u\in K$.

\begin{lemma}[Only ideals can be missed]\label{lem:missed-ideal}
The color set $H(T)$ ordinarily activated by a tree is upward closed.
Consequently $K(T)=\mathcal P\setminus H(T)$ belongs to
$\mathcal I(\mathcal P)$.  The same conclusion holds for the union of the
color masks of any subset of edge copies, whether or not that subset is
a spanning-tree basis.
\end{lemma}
\begin{proof}
If an edge row reaches $u$ and $u\preceq v$, transitivity implies that
it reaches $v$.  Each individual edge mask is therefore upward closed,
and so is a union of such masks.  Its complement is an order ideal.
Discarding free colors does not change this statement for the induced
partial order on the remaining colors.
\end{proof}

Retain the Boolean function $f$ from the preceding section, with its
free seeds included, but evaluate it only on the upward-closed masks
$\mathcal P\setminus K$.  Define
\[
                    g(K)=f(\mathcal P\setminus K)
                         \qquad(K\in\mathcal I(\mathcal P)).
\]
For an ideal $A$, let $\max(A)$ be its set of maximal elements and put
\begin{equation}\label{eq:ideal-coefficients}
                 c(A)=\sum_{S\subseteq\max(A)}
                             (-1)^{|S|}g(A\setminus S).
\end{equation}
Every set $A\setminus S$ occurring here is again an ideal.

\begin{lemma}[Explicit inversion on the ideal lattice]
\label{lem:ideal-inversion}
For every ideal $K$,
\begin{equation}\label{eq:ideal-inversion}
                 \sum_{A\in\mathcal I(\mathcal P),\,A\subseteq K}
                                   c(A)=g(K).
\end{equation}
Moreover the full coefficient table can be computed in $O(dJ)$ additions
and subtractions, without enumerating the subsets in
\eqref{eq:ideal-coefficients} separately.
\end{lemma}
\begin{proof}
Fix an ideal $B\subseteq K$ and collect the coefficient of $g(B)$ in the
left side of \eqref{eq:ideal-inversion}.  A contributing term is an ideal
$A=B\cup S$ with $S\subseteq\max(A)$.  This is equivalent to taking
$S$ to be an arbitrary subset of the minimal elements of $K\setminus B$.
Indeed, if $s\in A\setminus B$ is maximal in $A$, all its strict
predecessors lie in $B$: any predecessor outside $B$ would lie in
$A\setminus B$ and would not be maximal.  Hence $s$ is minimal outside
$B$.  Conversely, minimal elements outside an ideal $B$ are pairwise
incomparable and have every predecessor in $B$; adjoining any subset of
them gives an ideal in which those new elements are maximal.  Thus the
coefficient of $g(B)$ is
\[
             \sum_{S\subseteq\min(K\setminus B)}(-1)^{|S|}.
\]
It is one if $B=K$ and zero otherwise, since a nonempty finite poset
has a minimal element.  This proves the inversion identity over the
integers, and hence over every field.

For the efficient transform, choose a linear extension of the partial
order, with predecessors before successors.  Initialize $c(A)=g(A)$ for
all ideals.  Process colors $v$ in \emph{reverse} linear-extension order.
For every ideal $A$ for which $v\in A$ and $A\setminus\{v\}$ is an
ideal, perform
\begin{equation}\label{eq:restricted-transform}
                  c(A)\leftarrow c(A)-c(A\setminus\{v\}).
\end{equation}
The set-removal condition says exactly that $v$ is maximal in $A$.
To expand the coefficient eventually obtained for a fixed $A$, each
subtraction removes a maximal element and contributes one minus sign.
Any predecessor of a removed element occurs later in the processing order.
It has therefore not yet been processed in the source value when that
subtraction takes place.  Such a newly maximal predecessor cannot be
removed along the resulting expansion path.  Consequently the only
elements removed are members of the original $\max(A)$, and any subset
of them is removed exactly once, in reverse linear-extension order.
This gives exactly \eqref{eq:ideal-coefficients}.  There are at most $dJ$
updates.

The order matters.  For a chain $a\prec b$, reverse processing yields
$c(\{a,b\})=g(\{a,b\})-g(\{a\})$.  Processing predecessors first would
incorrectly introduce an extra $g(\varnothing)$ term in this coefficient.
\end{proof}

The transform is the M\"obius inversion for the distributive lattice of
order ideals; the argument above establishes the needed formula directly.
Efficient lattice zeta and inverse transforms are established tools
\cite{FastZeta2012,KaskiKohonenWesterback2016}.
No assumption about a new matroid of peelable sets is being made.

For each ideal $A$, let $P_A$ be the common-basis determinant that bans
copies whose color masks meet $A$, as in the preceding section.  Form
\begin{equation}\label{eq:ideal-filter}
                   Q_{\mathcal I}(z)
                       =\sum_{A\in\mathcal I(\mathcal P)}c(A)P_A(z).
\end{equation}

\begin{theorem}[Ideal-state determinant sieve]\label{thm:ideal-sieve}
The nonzero monomials of $Q_{\mathcal I}$ encode exactly the copied
spanning-tree bases leaving a completely peelable matrix.  In particular,
forest-compatible raw-completion witnesses can be found, with one-sided
miss probability at most $\varepsilon$, in
\begin{equation}\label{eq:ideal-time}
                     J\operatorname{poly}(N,\log(1/\varepsilon))
\end{equation}
time and $J\operatorname{poly}(N)$ space.  Every success is checked by
replaying a spanning tree and a unit-column peel order in the original
matrix.  The theorem applies to arbitrary exceptional supports through
their finite Horn rules.
\end{theorem}
\begin{proof}
Fix a copied common basis $F$, with underlying spanning tree $T$.  Its
monomial occurs in $P_A$ exactly when $A\subseteq K(T)$.  By
Lemma~\ref{lem:missed-ideal}, $K(T)$ is an ideal.  Therefore its coefficient
in \eqref{eq:ideal-filter} is its original nonzero basis coefficient
multiplied by
\[
       \sum_{A\in\mathcal I(\mathcal P),\,A\subseteq K(T)}c(A)
            =g(K(T))=f(H(T)).
\]
This is zero or one by Lemma~\ref{lem:ideal-inversion}.  The fixed source
groups are already enforced by the represented partition/filler matroid,
so the exact special-row criterion identifies precisely the successful
trees.  Different copied bases still have different monomials.  The
polynomial is homogeneous and multilinear of degree $r$.

Enumerate ideals by a depth-first traversal of a linear extension.
At the step for color $v$, exclusion is always allowed; inclusion is
allowed precisely when every predecessor of $v$ has already been included.
Each leaf is an ideal and every ideal occurs exactly once.  Every prefix
can be extended to a complete ideal by excluding all remaining elements,
so the number of prefixes at any fixed depth is at most $J$.  There are
at most $(d+1)J$ prefixes in total.  This gives
$J\operatorname{poly}(N)$ enumeration time, including exact predecessor
mask tests.

Compute $g$ on these $J$ states by the polynomial-time closure routine,
and compute $c$ by Lemma~\ref{lem:ideal-inversion}.  Each evaluation of
$Q_{\mathcal I}$ uses at most $J$ determinants and polynomial additional
work.  The same prime-field identity test and repeated-deletion witness
extraction as before apply, with unchanged degree and characteristic
requirements.  A polynomial number of determinant evaluations suffices
for extraction, and independent trials contribute only the stated
$\log(1/\varepsilon)$ factor.  Lists of ideal masks and their coefficient
values use $J\operatorname{poly}(N)$ space.
\end{proof}

\begin{corollary}[Width and chain bounds]\label{cor:ideal-width}
Suppose the implication poset admits a partition into $w\geq1$ chains
of lengths $n_1,\ldots,n_w$, where $\sum_i n_i=d$.  Then
\begin{equation}\label{eq:ideal-width}
       J\leq\prod_{i=1}^w(n_i+1)
               \leq\left(1+\frac d w\right)^w.
\end{equation}
In particular, if $w$ is the width of the poset, a chain partition of
that size exists by the classical finite chain-decomposition theorem of
Dilworth.  Fixed width gives randomized polynomial-time discovery for
this local completion problem even when $d$ is large.  If $N=n^{O(1)}$
and $w=O(\log n)$, the local time is quasi-polynomial in $n$.
When $d=0$, the unique ideal is the empty set and $J=1$.
\end{corollary}
\begin{proof}
A partition into $w$ chains exists by Dilworth's theorem
\cite{Dilworth1950}.  An ideal intersects each chain in a prefix, so it is uniquely determined
by the $w$ prefix lengths.  This injects the ideals into
$\prod_i\{0,\ldots,n_i\}$ and proves the first inequality.  The second
is the arithmetic--geometric mean inequality applied to the positive
numbers $n_i+1$.  Since $d\leq N$, the final complexity assertions follow
from Theorem~\ref{thm:ideal-sieve}.  A chain partition is needed only to
state or certify the bound; the ideal enumeration algorithm does not
need to find a minimum-width partition.
\end{proof}

This refinement changes the governing parameter from the number of
exceptional columns to the number of feasible ordinary activation patterns.
A chain has $d+1$ ideal states; an antichain has $2^d$ and recovers the full
Boolean bound.  Intermediate posets interpolate between these cases.
An antichain example shows why a small number of states cannot be assumed
merely from having compressed coordinates.  No bound on $J$, $d$, or $w$
for every state of an unknot hierarchy is established here.

\begin{theorem}[Deterministic discovery parameterized by implication width]
\label{thm:deterministic-width}
The same finite completion problem has a deterministic exact algorithm
with time
\[
                   J(m+1)^w\operatorname{poly}(N)
                            =N^{O(w+1)}
\]
and polynomial space, where $w$ is the width of the nonfree special-SCC
implication poset.  In particular fixed width gives deterministic
polynomial time, and $w=O(\log n)$ gives deterministic quasi-polynomial
local time when $N=n^{O(1)}$.
\end{theorem}
\begin{proof}
Stream the ideals $K$ and set $A=\mathcal P\setminus K$.  Discard the
state if $f(A)=0$.  Otherwise require a spanning tree to hit every fixed
mandatory source group and every set $\operatorname{Anc}(v)$ for
$v\in\min(A)$.  The set of minima is an antichain and has at most $w$
elements.  Enumerate one eligible edge representative for each such
ancestor requirement.  Reject a selection containing a cycle; contract
its forest; discard the mandatory groups it has already hit; and use
graphic--partition matroid intersection for the remaining disjoint groups,
exactly as in Theorem~\ref{thm:triple-xp}.  Extend a successful forest to
a spanning tree and verify the original peel order.

For soundness, an upward-closed set containing every minimal element of
$A$ contains all of $A$.  Thus the tree's ordinary influence satisfies
$A\subseteq H(T)$, and monotonicity gives $f(H(T))\geq f(A)=1$.
The fixed mandatory groups complete the exact special-row criterion.
For completeness, any successful tree occurs in the case $A=H(T)$;
it has a representative in each of the corresponding ancestor sets,
and those representatives are contained in the tree and hence independent.
The remaining group representatives can also be selected from that tree.

There are $J$ states and at most $(m+1)^w$ representative tuples per state.
The width bound on $J$ proves the stated input-size bound.  Depth-first
ideal enumeration, a single closure computation, and a single
representative choice can all be retained at a time.  The matroid
intersection computation uses polynomial space.  The zero-dimensional
case has one state and no ancestor requirements.
\end{proof}
